\documentclass[11pt]{article}
\usepackage[a4paper,margin=1in]{geometry}
\usepackage{amsmath,amssymb,amsthm}
\usepackage[T1]{fontenc}
\usepackage{lmodern}
\usepackage[hidelinks]{hyperref}
\setlength{\emergencystretch}{3em}

\theoremstyle{plain}
\newtheorem{theorem}{Theorem}
\newtheorem{lemma}[theorem]{Lemma}
\newtheorem{proposition}[theorem]{Proposition}
\newtheorem{corollary}[theorem]{Corollary}
\theoremstyle{definition}
\newtheorem{definition}[theorem]{Definition}
\theoremstyle{remark}
\newtheorem{remark}[theorem]{Remark}

\DeclareMathOperator{\per}{per}

\title{A Proof of Conjecture 00000000050}
\author{%
  Feiyu\thanks{Independent researcher.}
  \and
  Claude (Opus 5)\thanks{Anthropic.}%
}
\date{2026-09-12}

\begin{document}
\maketitle

\begin{abstract}
Conjecture 00000000050 asserts that every prime is the permanent of a square
matrix with entries in $\{0,1\}$. We prove it, and more: every natural number
is, by an explicit $(n+1)\times(n+1)$ construction whose permanent is $n+1$.
Primality plays no role. The proof is formalized in Lean~4 against Mathlib,
using Mathlib's \texttt{Matrix.permanent}.
\end{abstract}

\section{Statement}

For a square matrix $A=(a_{rc})$ over a commutative semiring, indexed by a
finite set $I$, the permanent is
\[
  \per A=\sum_{\sigma\in\mathfrak S_I}\ \prod_{i\in I}a_{\sigma(i),\,i} ,
\]
the determinant without signs. Conjecture 00000000050 asserts that for every
prime $p$ there is a square matrix with entries in $\{0,1\}$ whose permanent is
exactly $p$.

\section{The construction}

\begin{definition}
For $n\ge0$ let $W_n$ be the $(n+1)\times(n+1)$ matrix with rows and columns
indexed by $\{0,1,\ldots,n\}$ and entries
\[
  (W_n)_{r,c}=
  \begin{cases}
    1,&\text{if }r=0\text{ or }c=0\text{ or }c=r,\\
    0,&\text{otherwise.}
  \end{cases}
\]
\end{definition}

That is: the first row is all ones, and row $r\ne0$ has ones exactly in
column $0$ and column $r$. For $n=4$,
\[
  W_4=
  \begin{pmatrix}
    1&1&1&1&1\\
    1&1&0&0&0\\
    1&0&1&0&0\\
    1&0&0&1&0\\
    1&0&0&0&1
  \end{pmatrix}.
\]
Its entries lie in $\{0,1\}$ by construction.

\begin{lemma}\label{lem:contrib}
A permutation $\sigma$ of $\{0,\ldots,n\}$ satisfies
$\prod_i (W_n)_{\sigma(i),i}=1$ if and only if
\begin{equation}\label{eq:contrib}
  \sigma(i)\in\{0,i\}\qquad\text{for every }i\ne0 ,
\end{equation}
and $\prod_i (W_n)_{\sigma(i),i}=0$ otherwise.
\end{lemma}

\begin{proof}
All entries are $0$ or $1$, so the product is $1$ exactly when every factor is,
and $0$ otherwise. The factor at index $i$ is $(W_n)_{\sigma(i),i}$, which is
$1$ precisely when $\sigma(i)=0$ or $i=0$ or $i=\sigma(i)$. Requiring this for
every $i$ is exactly \eqref{eq:contrib}, the index $i=0$ imposing nothing.
\end{proof}

\begin{lemma}\label{lem:swaps}
The permutations satisfying \eqref{eq:contrib} are exactly the transpositions
$\tau_k=(0\ k)$ for $k\in\{0,\ldots,n\}$, where $\tau_0$ is the identity. They
are pairwise distinct, so there are $n+1$ of them.
\end{lemma}

\begin{proof}
Each $\tau_k$ satisfies \eqref{eq:contrib}: for $i\ne0$ we have $\tau_k(i)=0$
if $i=k$, and $\tau_k(i)=i$ otherwise.

Conversely let $\sigma$ satisfy \eqref{eq:contrib} and put $k=\sigma(0)$; we
claim $\sigma=\tau_k$. At $0$ both sides give $k$. Let $i\ne0$.

If $\sigma(i)=i$, then $i\ne k$: otherwise $\sigma(0)=k=i=\sigma(i)$ would give
$0=i$. Hence $\tau_k(i)=i=\sigma(i)$.

If instead $\sigma(i)=0$, we show $i=k$. First $k\ne0$, since $k=0$ would give
$\sigma(0)=0=\sigma(i)$ and so $i=0$. Applying \eqref{eq:contrib} at $k\ne0$
gives $\sigma(k)\in\{0,k\}$; and $\sigma(k)=k$ is impossible, since with
$\sigma(0)=k$ it would force $k=0$. So $\sigma(k)=0=\sigma(i)$, whence $i=k$
and $\tau_k(i)=\tau_k(k)=0=\sigma(i)$.

Finally $k\mapsto\tau_k$ is injective because $\tau_k(0)=k$.
\end{proof}

\begin{theorem}\label{thm:per}
$\per W_n=n+1$ for every $n\ge0$.
\end{theorem}

\begin{proof}
By Lemma~\ref{lem:contrib} the sum defining $\per W_n$ counts the permutations
satisfying \eqref{eq:contrib}, each contributing $1$; by Lemma~\ref{lem:swaps}
there are $n+1$ of them.
\end{proof}

\section{The conjecture}

\begin{corollary}\label{cor:all}
Every natural number $m$ is the permanent of a square matrix with entries in
$\{0,1\}$.
\end{corollary}

\begin{proof}
For $m=0$ take the $1\times1$ zero matrix. For $m=k+1$ take $W_k$ and apply
Theorem~\ref{thm:per}.
\end{proof}

\begin{corollary}
Conjecture 00000000050 is true.
\end{corollary}

\begin{proof}
A prime is a natural number; apply Corollary~\ref{cor:all}. Explicitly, for a
prime $p$ the matrix $W_{p-1}$ is a $p\times p$ matrix with entries in
$\{0,1\}$ and $\per W_{p-1}=p$.
\end{proof}

\begin{remark}
The hypothesis that $p$ be prime is not used anywhere, and cannot be: the
permanents of $\{0,1\}$ matrices are exactly the natural numbers, by
Corollary~\ref{cor:all} in one direction and because a permanent of a
nonnegative integer matrix is a nonnegative integer in the other. A
conjecture restricted to primes is therefore a special case of a statement with
no arithmetic content. Read as a question about the \emph{smallest} matrix with
a given permanent --- the minimal $N$ for which some $N\times N$ $\{0,1\}$
matrix has permanent $p$ --- it would not be vacuous; our construction gives
$N\le p$, and we do not address minimality.
\end{remark}

\begin{remark}
The construction has a transparent combinatorial reading: $\per A$ counts the
perfect matchings of the bipartite graph with biadjacency matrix $A$. For
$W_n$, rows $1,\ldots,n$ each see only columns $0$ and their own, so a perfect
matching either matches every $r\ne0$ to column $r$ --- leaving row $0$ for
column $0$ --- or matches a single $r\ne0$ to column $0$, forcing row $0$ to
column $r$. That is $n+1$ matchings.
\end{remark}

\section{Formal verification}

The accompanying Lean~4 project formalizes the proof against Mathlib. The
conjecture is stated in the formalization, so the final theorem is a
formalization of the conjecture itself. Permanents are Mathlib's
\texttt{Matrix.permanent}, not a local redefinition.

\begin{itemize}
  \item \texttt{W}, \texttt{W\_apply}, \texttt{W\_zero\_or\_one} --- the matrix
    and the fact that its entries lie in $\{0,1\}$.
  \item \texttt{Contributes}, \texttt{prod\_W} --- Lemma~\ref{lem:contrib}: the
    product over a permutation is $1$ or $0$ according to \eqref{eq:contrib}.
  \item \texttt{contributes\_swap}, \texttt{eq\_swap\_of\_contributes},
    \texttt{swap\_injective}, \texttt{filter\_contributes} ---
    Lemma~\ref{lem:swaps}: the contributing permutations are exactly the
    $n+1$ transpositions $(0\ k)$.
  \item \texttt{permanent\_W} --- Theorem~\ref{thm:per}.
  \item \texttt{exists\_zeroOne\_permanent} --- Corollary~\ref{cor:all}.
  \item \texttt{ConjectureHolds}, \texttt{conjecture\_00000000050\_true} ---
    the conjecture and its proof.
\end{itemize}

The project builds with \texttt{lake build} on
\texttt{leanprover/lean4:v4.33.1} with Mathlib \texttt{v4.33.1}. It contains no
\texttt{sorry}, no \texttt{admit} and no \texttt{native\_decide}, and
introduces no axioms of its own:
\texttt{\#print axioms conjecture\_00000000050\_true} reports exactly the three
standard axioms \texttt{propext}, \texttt{Classical.choice} and
\texttt{Quot.sound}.

\end{document}
